\section{Related Work}

Delay-tolerant networking was introduced to support communication in environments where conventional Internet assumptions break down, particularly when end-to-end paths are intermittent or absent for long periods~\cite{fall2003dtn}. In these networks, store-carry-forward communication replaces route-then-send behavior. This shift makes forwarding strategy the central protocol problem: each encounter is a rare opportunity, and routing decisions must be made with partial, local information. Several surveys and overviews have since organized DTN routing into broad families based on the information available to the forwarding decision, the amount of replication allowed, and whether the protocol uses mobility history, social structure, or explicit optimization goals~\cite{zhang2006overview,sobin2016survey}.

One early and influential baseline is Epidemic Routing, which forwards data aggressively to maximize delivery probability~\cite{vahdat2000epidemic}. Its main weakness is that it consumes substantial bandwidth, storage, and energy. Jain et al. helped frame DTN routing as a distinct problem class in which different algorithms can rely on different levels of connectivity knowledge~\cite{jain2004routing}. Later protocols explored more structured forwarding choices. PRoPHET uses encounter history and transitivity to prefer relays with stronger delivery likelihoods~\cite{lindgren2012prophet}, while MaxProp prioritizes transmission scheduling and path likelihoods in highly disrupted vehicular networks~\cite{burgess2006maxprop}. RAPID takes a different angle by formulating DTN forwarding as a resource-allocation problem and explicitly optimizing routing metrics such as average delay~\cite{balasubramanian2007rapid}. Delegation Forwarding further shows that useful forwarding policies can emerge from threshold-style relay selection without relying on unrestricted flooding~\cite{erramilli2008delegation}.

Another important branch of work uses social structure rather than only raw contact frequency. Bubble Rap, for example, exploits community structure and node centrality to guide forwarding toward more socially informative relays~\cite{hui2011bubblerap}. These social-aware protocols are particularly effective when encounter opportunities are shaped by recurring human mobility patterns rather than purely random contacts. They are conceptually richer than the protocol studied here, but they also require substantially more accumulated context than a bounded-replication method such as Spray and Wait.

Spray and Wait was proposed to preserve much of epidemic routing's delivery benefit while limiting replication through a strict copy budget~\cite{spyropoulos2005spraywait}. The later multiple-copy analysis further formalized why bounded replication can achieve strong performance in intermittently connected settings~\cite{spyropoulos2008multiplecopy}. Relative to PRoPHET, MaxProp, RAPID, Delegation Forwarding, and Bubble Rap, Binary Spray and Wait is attractive for this project because it isolates the replication tradeoff in a particularly transparent way: the key control parameter is simply the number of injected copies. Subsequent evaluations of opportunistic routing have emphasized that delivery ratio, latency, and transmission cost must be considered together rather than in isolation~\cite{song2007evaluating,sobin2016survey}.

This project fits into that line of work as a small, simulator-based reproduction of the central Spray and Wait idea. Its novelty is not the invention of a new DTN routing protocol, but a careful implementation and evaluation of Binary Spray and Wait inside a minimal wireless teaching simulator. That scope is appropriate for the project: it is original relative to the provided PyWiSim examples, technically meaningful, and small enough to study transparently end to end.
